\section{Scope and method}\label{sec:method}

This paper describes a proof that exists as Lean source. It was written by reading that source
from the last line upward: each hypothesis was followed to the declaration that supplies it,
and each definition was transcribed into ordinary notation, until every branch reached a
theorem of the published literature. Three rules governed the reading.

\emph{Definitions are transcribed, not inferred from names.} Names in the source are labels.
The same name is used for different propositions in different namespaces, and several
formulations of one step coexist.

\emph{Membership in the proof is decided by the compiled environment.} Almost every file of the
archive imports almost every other, so the text alone cannot show that a declaration is used by
the final theorem. The dependency cone of \cref{sec:cone} decides it by machine, and every
statement in this paper that something is or is not used by the proof refers to it.

\emph{Every statement about the proof has a recorded basis.} \Cref{app:audit} lists the
statements of this paper with the file and line where each was read.

The paper does not reprove anything, and it describes what the proof does, not how it was
found. Constants and lower-order terms are abbreviated; each inequality displayed here holds in
the source with explicit constants. The proofs of the published theorems that the archive
formalises were not read; \cref{sec:literature} identifies them.

Qui Vincit directed the work and set its standards. Fable, a language model (Claude,
Anthropic), ran the build, read the source and wrote the text, with further instances of the
model reading bounded parts of the source; their reports are published with the paper.
Corrections are welcome at the repository of \cref{app:reproduce}.
